% Simplified teaching-learning plan.
% Compile with pdfLaTeX.
\documentclass[10pt,letterpaper]{article}

\usepackage[margin=0.6in,bottom=0.72in]{geometry}
\usepackage[T1]{fontenc}
\usepackage[utf8]{inputenc}
\usepackage[scaled=0.95]{helvet}
\renewcommand{\familydefault}{\sfdefault}
\usepackage{microtype}
\usepackage{enumitem}
\usepackage{fancyhdr}
\usepackage{needspace}

\setlength{\parindent}{0pt}
\setlength{\parskip}{3pt}

\pagestyle{fancy}
\fancyhf{}
\renewcommand{\headrulewidth}{0pt}
\renewcommand{\footrulewidth}{0pt}
\fancyfoot[L]{\small Versión 08}
\fancyfoot[R]{\small\thepage}

\newcommand{\criterion}[2]{\item[\textbf{C#1.}] #2}
\newcommand{\labelled}[2]{\textbf{#1:} #2\par}
\newcommand{\plansection}[1]{\par\vspace{8pt}{\large\bfseries #1\par}\vspace{3pt}}
\newcommand{\plansubsection}[1]{\Needspace{6\baselineskip}\par\vspace{4pt}{\normalsize\bfseries #1\par}\vspace{1pt}}

\begin{document}

\vspace*{-0.25in}
\begin{center}
  {\large\bfseries TEACHING-LEARNING PLAN}\par
  \vspace{2pt}
  \textbf{PGF-02-R39}\par
  \textbf{Date:} August 12 -- November 5, 2026
\end{center}

\plansection{1. Identification}

\labelled{Area of knowledge}{Global Economics}
\labelled{Teacher's name (or Teachers' names)}{Nicolas Lopez}
\labelled{Grade}{10}
\labelled{Term}{1 (2026-2027)}

\plansection{2. Learning Objective}

\plansubsection{Learning Objective}

Analyze economic decisions under uncertainty by formulating alternatives,
events, states and consequences; constructing payoff tables from direct payoffs
or revenue and cost information; applying and comparing decision criteria with
and without probabilities; and communicating a justified recommendation
through tables, decision trees and project evidence.

\plansubsection{Assessment Criteria}

\begin{description}[leftmargin=3.1em,labelwidth=2.7em,labelsep=0.4em,itemsep=2pt,topsep=2pt]
  \criterion{1}{Characterize a decision under uncertainty by identifying its context, alternatives, event, states of nature and consequences.}
  \criterion{2}{Interpret and organize alternatives, events, states, consequences and supplied payoffs in a payoff table.}
  \criterion{3}{Calculate revenue, variable and fixed costs, quantities and profit, and construct the corresponding payoff table.}
  \criterion{4}{Apply Maximax and Maximin to select and justify optimistic and conservative alternatives without probabilities.}
  \criterion{5}{Construct opportunity-loss tables and apply Minimax Regret.}
  \criterion{6}{Use probabilities to calculate expected monetary value and select the preferred alternative.}
  \criterion{7}{Analyze the sensitivity of expected-value recommendations to changes in probabilities.}
  \criterion{8}{Compare Bayes/expected value, Maximin and Minimax Regret on common decision problems.}
  \criterion{9}{Represent staged decisions and uncertainty with decision trees.}
  \criterion{10}{Conduct and communicate an end-to-end decision analysis that integrates models, criteria, results and a justified recommendation.}
\end{description}

\plansection{3. Conceptual standards of the disciplinary knowledge that supports the narrative's development}

\textbf{Core disciplinary knowledge:} decisions under uncertainty;
decision-makers; alternatives; events; states of nature; consequences; direct
and calculated payoffs; revenue, quantities, variable and fixed costs, profit
and payoff tables.

\textbf{Analytical standards:} Maximax, Maximin, opportunity loss and Minimax
Regret; probabilities, Bayes/expected monetary value and sensitivity analysis;
decision trees; integrated comparison, interpretation and justification of
recommendations.

\plansection{4. Pre-lesson and Context}

\textit{Pedagogical actions planned by the teacher to motivate students for the
term work proposal (presentation of scenarios, narratives, learning objectives
and assessment criteria).}

Begin with familiar uncertain situations from retail, transport, energy and
school operations. Use the Term 1 survey and a short initial scenario to elicit
students' decision processes, then introduce the term learning objective,
C1-C10 criteria, the three learning-evidence blocks and the payoff-table
structure. Share materials and instructions through OPEN/Teams.

\plansection{5. Experience}

\textit{Pedagogical actions that establish the learning objective according to
the assessment criteria. Links to virtual resources or ways used to share
experiences (OPEN, TEAMS, etc.) should be included.}

\plansubsection{Pedagogical Action N1 -- C1-C4}

\textit{First Learning Evidence}

\labelled{Learning sequence}{Develop Practice Activities 1-3 through teacher
modeling, guided analysis and individual or paired work. Students identify
decision elements, organize and calculate payoffs, construct payoff tables, and
apply Maximax and Maximin.}

\labelled{Learning evidence}{Written exam based on the practice activities,
assessing C1-C4. Use the criterion-specific catch-up materials for feedback and
remediation.}

\plansubsection{Pedagogical Action N2 -- C5-C7}

\textit{Second Learning Evidence}

\labelled{Learning sequence}{Complete the regret component of Practice Activity
3 and Practice Activity 4. Students construct opportunity-loss tables, apply
Minimax Regret, calculate expected monetary value, and test how recommendations
change as probabilities vary.}

\labelled{Learning evidence}{Written exam based on the practice activities,
assessing C5-C7. Follow with targeted correction and criterion-specific
catch-up work when needed.}

\plansubsection{Pedagogical Action N3 -- C8-C10}

\textit{Third Learning Evidence}

\labelled{Learning sequence}{Use the C8-C10 extension practice to compare
Bayes/expected value, Maximin and Minimax Regret on common data; construct
decision trees; and complete an end-to-end analysis of an economic decision
under uncertainty.}

\labelled{Learning evidence}{Written exam based on the practice activities,
complemented by project work in which students model a case, apply the relevant
methods, interpret results and justify a final recommendation.}

\plansection{References}

\begin{itemize}[leftmargin=1.4em,itemsep=2pt,topsep=2pt]
  \item \textbf{Course map:} README.md -- Grade 10 Global Economics, Term 1.
  \item \textbf{Practice Activity 1:} \textnormal{\detokenize{G10_T1_L1_C1C2_practice_activity1.tex}} -- decision elements and payoff tables.
  \item \textbf{Practice Activity 2:} \textnormal{\detokenize{G10_T1_L1_C1C3_practice_activity2.tex}} -- revenue, costs, profit and payoff tables.
  \item \textbf{Practice Activity 3:} \textnormal{\detokenize{G10_T1_L2_C4C5_practice_activity3.tex}} -- Maximax, Maximin and Minimax Regret.
  \item \textbf{Practice Activity 4:} \textnormal{\detokenize{G10_T1_L2_C6C7_practice_activity4.tex}} -- expected value and sensitivity analysis.
  \item \textbf{Extension material:} \textnormal{\detokenize{G10_T1_L3_MATERIAL.tex}} -- method comparison, decision trees and integrated C8-C10 analysis.
\end{itemize}

\end{document}
